% =============================================================
% chapters/rozdil3.tex
% РОЗДІЛ 3. Веб-розробка інформаційно-аналітичної системи
% на основі архітектури «ПрофОсвіта Луганщини» v2.0.
% Цільовий обсяг: ~17 стор. (Phase D плану).
% =============================================================

\section*{\centering РОЗДІЛ 3.\\ ВЕБ-РОЗРОБКА ІНФОРМАЦІЙНО-АНАЛІТИЧНОЇ СИСТЕМИ}
\addcontentsline{toc}{section}{РОЗДІЛ 3. ВЕБ-РОЗРОБКА ІНФОРМАЦІЙНО-АНАЛІТИЧНОЇ СИСТЕМИ}
\label{sec:rozdil3}

% TODO Phase D — наповнити повністю.
% Джерела (RO):
%   project/docs/superpowers/plans/Стуктура вебсайту/Архітектура_сайту_ПрофОсвіта_Луганщини_v2.docx
%   project/docs/superpowers/plans/Стуктура вебсайту/Specifikatsiya_24_form_Monitoringu_v2.md
%   project/docs/superpowers/plans/Стуктура вебсайту/struktura_saytu_NMC_PO.md
%   14 HTML-сторінок прототипу (index, dashboard, institutions, monitoring*, …)
% Скриншоти: bachelor_thesis/media/screenshots/*.pdf (8–12 шт., згенерувати у Phase D)

\subsection*{3.1. Технологічний стек і засоби розробки}
\addcontentsline{toc}{subsection}{3.1. Технологічний стек і засоби розробки}

% TODO Phase D: 2 стор. — HTML5/CSS3/JS/Tailwind/Chart.js (прототип) +
% Next.js 14/TypeScript/PostgreSQL (виробничий стек). Обґрунтування проти Vue/Django.

\subsection*{3.2. Реалізація публічної частини порталу}
\addcontentsline{toc}{subsection}{3.2. Реалізація публічної частини порталу}

% TODO Phase D: 4 стор. — Головна, Дашборд, Заклади освіти, Видання, Про проєкт,
% Контакти. Скриншоти: Рис. 3.1–3.4 із bachelor_thesis/media/screenshots/.

\subsection*{3.3. Реалізація закритої частини моніторингу}
\addcontentsline{toc}{subsection}{3.3. Реалізація закритої частини моніторингу}

% TODO Phase D: 4 стор. — авторизація, кабінет ЗПО, 24 форми за 7 блоками,
% перехресна валідація, адмін-панель НМЦ. Скриншоти: Рис. 3.5–3.6.
% Табл. 3.1: 24 форми моніторингу за 7 блоками (з Specifikatsiya_24_form_Monitoringu_v2.md).

\subsection*{3.4. Реалізація візуалізації даних}
\addcontentsline{toc}{subsection}{3.4. Реалізація візуалізації даних}

% TODO Phase D: 2 стор. — мапа релокації (Leaflet + OSM), KPI-дашборд (Chart.js),
% графік динаміки контингенту 2022–2026 із вертикальним маркером 24.02.2022,
% експорт у Excel (SheetJS) і PDF (jsPDF).

\subsection*{3.5. Дизайн-система, адаптивність, доступність}
\addcontentsline{toc}{subsection}{3.5. Дизайн-система, адаптивність, доступність}

% TODO Phase D: 2 стор. — палітра 6 кольорів (бренд #185FA5 та ін.),
% типографіка Inter + Lora, breakpoints (desktop ≥1024px, tablet 768–1023px,
% mobile ≤767px), WCAG 2.1 AA contrast ≥4.5:1, alt-texts, навігація з клавіатури.

\subsection*{3.6. Результати апробації прототипу}
\addcontentsline{toc}{subsection}{3.6. Результати апробації прототипу}

% TODO Phase D: 2 стор. — експертне оцінювання (≥3 експерти: керівник,
% рецензент, директор НМЦ ПТО Луганщини). Методологія: SUS + домен-специфічні
% пункти. Результати: юзабіліті, ергономіка, функціональна повнота.
% Виявлені недоліки + план виправлень. Дорожня карта Q2 2026 – Q1 2027.

\subsection*{Висновки до розділу 3}

% TODO Phase D: 1 стор., 4–5 пунктів. На що відповідає Розділ 3:
% (1) реалізовано публічну частину з мапою релокації та KPI-дашбордом;
% (2) реалізовано закриту частину з 24 формами моніторингу;
% (3) забезпечено WCAG 2.1 AA доступність;
% (4) проведено експертне оцінювання (TRL 4–5);
% (5) визначено дорожню карту до пілотного впровадження.

\clearpage
